\chapter{Core Mechanics}
\label{ch:core-mechanics}

\begin{tcolorbox}[colback=black!3!white,colframe=black!40!white,title=\textit{A Word from the Raven Road},breakable]
\emph{``The road remembers. Every broken wheel leaves a mark, every lit lamp bears witness. The only question is: what are you willing to owe?''}  
\vspace{0.2cm}\hfill --- Dusana of the Raven Road, \textit{The Hearth Ledger}
\end{tcolorbox}

\section{The 90-Second Play Loop}
\label{sec:play-loop}
\index{play loop}\index{GM loop}

At the heart of every scene is a simple, fast loop that keeps the story moving:

\begin{enumerate}
  \item \textbf{GM sets the stage} -- Describe the situation and ask ``What do you do?''
  \item \textbf{Player declares intent and approach} -- State what you want to achieve and how (Attribute + Skill).
  \item \textbf{GM sets Position and Difficulty Value (DV)} -- Choose Dominant/Controlled/Desperate and a DV (2--5+).
  \item \textbf{Player rolls dice} -- Count Successes (6+) and Story Beats (1s).
  \item \textbf{Apply outcome} -- Compare Successes to DV (Clean Success / Success with SB / Partial / Miss).
  \item \textbf{GM spends Story Beats (SB)} -- Introduce complications, advance timers, or escalate the scene.
  \item \textbf{Player marks resources} -- Gain Boons (1 on Partial, 2 on Miss), track Fatigue/Harm, update Obligation, etc.
  \item \textbf{Pass the spotlight} -- Next player acts, or GM frames a new beat.
\end{enumerate}

\noindent This loop should take about 90 seconds per player action. If it drags, simplify: focus on intent, one roll, and a clear consequence.

\section{The Core Resolution Cycle}
\label{sec:core-resolution-cycle}
\index{resolution cycle}

Every meaningful action in Fate's Edge follows this cycle:

\begin{enumerate}
  \item \textbf{Declare Action \& Approach:} State what you want to achieve and how (Attribute + Skill).
  \item \textbf{GM sets Difficulty Value (DV):} Based on narrative stakes (see \ref{sec:difficulty-value}).
  \item \textbf{GM sets Position:} Dominant, Controlled, or Desperate (see \ref{sec:position}).
  \item \textbf{Roll dice pool:} Attribute + Skill d10s.
  \item \textbf{Count Successes (6+) and Story Beats (1s).}
  \item \textbf{Check against DV using the Outcome Matrix (see \ref{sec:outcome-matrix}).}
  \item \textbf{Apply outcome:} Clean Success, Success with SB, Partial, or Miss.
  \item \textbf{GM spends Story Beats (SB)} for complications and twists.
\end{enumerate}

\begin{tcolorbox}[colback=yellow!5!white,colframe=yellow!75!black,title=\textbf{Golden Rule},breakable]
\index{Golden Rule}
\textbf{When in doubt, make the choice that serves the story.} If you find yourself asking ``What's the optimal move?'' try reframing it as ``What would be most interesting for the story?''
\end{tcolorbox}

\section{Player-Managed Modules}
\label{sec:player-managed-modules}
\index{player trackers}\index{resources!player-managed}

In Fate's Edge, you are the primary manager of your character's mechanical state. This isn't just about convenience---it is a core design principle that reduces GM overhead and keeps you engaged.

\begin{tcolorbox}[colback=black!3!white,colframe=black!40!white,title=\textbf{You Track Your Own Story},breakable]
\textbf{What You Manage:}
\begin{itemize}
    \item Obligation tracks (one per Patron or Symbol)
    \item Corruption timers (Cantors, Pact-Whisperers)
    \item Leash timers (summoned spirits)
    \item Asset states (Flourishing, Stable, Strained, Collapsed)
    \item Follower states (Loyalty: Faithful/Strained/Broken; Fitness: Ready/Hurt/Down)
    \item Progression timers (Repertoire, Spirit Bond)
    \item Boons, Favor, Leverage
\end{itemize}

\textbf{Your Responsibilities:}
\begin{enumerate}
    \item \textbf{Mark immediately} when a rule says ``mark +X''.
    \item \textbf{Declare thresholds} (``My Obligation is full!'').
    \item \textbf{Keep it visible} (use trackers everyone can see).
    \item \textbf{Maintain ownership} -- you know your state best.
\end{enumerate}

\textbf{GM Support:} The GM spot-checks and provides consequences, but you drive your character's mechanical narrative.
\end{tcolorbox}

\begin{tcolorbox}[colback=black!3!white,colframe=black!40!white,title=\textbf{Example: Tracking Your Obligation},breakable]
  A character with Spirit 2 and Presence 2 has an Obligation Capacity of 4.  
  She invokes a Rite from Patron A (+1), then later Pushes a Rite from Patron B (+1), then Cracks the Seal on a Symbol (+2).  
  She now has:
  \begin{itemize}
    \item Patron A: 1 segment
    \item Patron B: 1 segment
    \item Symbol (Invoker): 2 segments
  \end{itemize}
  Total Obligation = 4 $\rightarrow$ exactly at Capacity. No Fatigue yet.  
  If she invokes one more Rite from any Patron, she exceeds Capacity and marks \textbf{1 Fatigue} immediately.
\end{tcolorbox}

\section{Basic Dice Mechanics}
\label{sec:dice-mechanics}
\index{dice}\index{Success}\index{Story Beats}

When you attempt a significant action, roll a pool of ten-sided dice (d10s). The size of your pool is:

\[
\text{Dice Pool} = \text{Attribute} + \text{Skill}
\]

\begin{description}
  \item[Attribute] (1--5) Broad traits like Strength, Wit, or Charm.
  \item[Skill] (0--5) Training in a specific area (e.g., Melee, Stealth, Arcana).
\end{description}

\begin{tcolorbox}[colback=red!5!white,colframe=red!75!black,title=\textit{Lore of the Fall: The Weight of Numbers},breakable]
\emph{``The Aeler count their breaths; the Aelaerem pour eight cups and stop. The ninth is not a number -- it is a wound that hasn't closed.''}  
In Fate's Edge, dice are not mere randomisers. They are the echo of every oath sworn, every debt unpaid. When you roll a 1, you feed the Hollow -- and the Hollow feeds the story.
\end{tcolorbox}

\subsection*{Reading the Dice}
\index{Success}\index{Story Beats}

Each die that rolls \textbf{6 or higher} counts as a \textbf{Success}.  
Each die that rolls a \textbf{1} generates a \textbf{Story Beat (SB)}.

\begin{center}
\begin{tabular}{lc}
\toprule
\textbf{Die Result} & \textbf{Effect} \\
\midrule
6--10 & +1 Success \\
1 & +1 Story Beat (SB) \\
2--5 & Nothing \\
\bottomrule
\end{tabular}
\end{center}

\paragraph{Example:} Lyra (Agility 3 + Stealth 2 = 5 dice) rolls: 6, 4, 3, 1, 6 $\rightarrow$ 2 Successes, 1 SB. The DV is 2 -- she succeeds, but the GM has 1 SB to spend (perhaps a guard stirs).

\begin{tcolorbox}[colback=purple!5!white,colframe=purple!75!black,title=\textbf{When to Roll},breakable]
  You roll only when:
  \begin{itemize}
    \item The outcome is uncertain, \textbf{and}
    \item Failure would matter to the story, \textbf{and}
    \item The action is significant enough to deserve a spotlight.
  \end{itemize}
  For routine actions with no real risk, the group can simply agree that you succeed.  
  \textit{Example:} Climbing a knotted rope in calm weather $\rightarrow$ no roll.  
  Climbing the same rope while goblins shoot arrows $\rightarrow$ roll.
\end{tcolorbox}

\section{Difficulty Value (DV)}
\label{sec:difficulty-value}
\index{Difficulty Value (DV)}

The Game Master sets a \textbf{Difficulty Value} -- the target number of Successes needed.

\begin{tcolorbox}[colback=gray!5!white,colframe=gray!70!black,title=\textbf{Difficulty Value (DV) Ladder},breakable]
\begin{center}
\begin{tabular}{cl}
\toprule
\textbf{DV} & \textbf{Situation} \\
\midrule
2 & Routine, no pressure \\
\textbf{3} & \textbf{Pressured, mild opposition (default)} \\
4 & Difficult, active resistance \\
5 & Extreme, high stakes \\
6+ & Nearly impossible, mythic challenge \\
\bottomrule
\end{tabular}
\end{center}
\end{tcolorbox}

\noindent \textbf{DV 3 is the default.} If the GM calls for a roll, something is at stake. Do not inflate DVs -- they already do their job.

\section{The Outcome Matrix}
\label{sec:outcome-matrix}
\index{Outcome Matrix}\index{Partial success}\index{Miss}

Compare your Successes to the DV to determine the outcome.

\begin{tcolorbox}[colback=gray!5!white,colframe=gray!70!black,title=\textbf{Outcome Matrix},breakable]
\begin{center}
\begin{tabular}{ll}
\toprule
\textbf{Outcome} & \textbf{What It Means} \\
\midrule
\textbf{Clean Success} (S $\geq$ DV) & You achieve your goal without complication. \\
\textbf{Success with SB} (S $\geq$ DV, SB>0) & You succeed, but the GM spends SB to add a twist. \\
\textbf{Partial} (0 < S < DV) & You make progress, but the situation remains unresolved. You gain 1 Boon. \\
\textbf{Miss} (S = 0) & You fail, and the situation escalates. You gain 2 Boons. \\
\bottomrule
\end{tabular}
\end{center}
\end{tcolorbox}

\paragraph{What a Partial Means.}
A Partial is not failure -- it moves the story forward. Examples:
\begin{itemize}
  \item Improve your Position for the next attempt (see \ref{sec:position}).
  \item Reduce the effective DV of the next roll.
  \item Succeed at reduced scale or duration.
  \item Reveal information or open a new path.
\end{itemize}

\begin{tcolorbox}[colback=red!5!white,colframe=red!75!black,title=\textbf{Fail Forward},breakable]
  A \textbf{Miss} is never ``nothing happens.''  
  The GM must introduce a complication, advance a timer, or reveal a hidden danger.  
  You gain 2 Boons -- the story moves, and you get resources to try again.
  
  \textit{Example:} You fail to pick a lock (Miss). The GM says: ``Your pick snaps, but you hear the guard's footsteps -- he's about to turn the corner. You have one beat to hide.'' The story advances.
\end{tcolorbox}

\subsection{Critical Successes}
\label{sec:critical-success}
\index{Critical success}\index{10s}

A roll of \textbf{10} counts as \textbf{two Successes}. If the roll succeeds, each 10 adds narrative weight.

\begin{center}
\begin{tabular}{ll}
\toprule
\textbf{10s} & \textbf{Effect} \\
\midrule
1 & Strong success -- improve your Position by one step on your next action this scene, or gain a minor advantage (GM's discretion). \\
2 & Exceptional success -- choose two benefits from the list above. \\
3 & Legendary success -- decisive victory; clear 1 segment from a relevant timer. \\
4+ & Mythic success -- reshape the scene; clear 1--2 segments from a relevant timer. \\
\bottomrule
\end{tabular}
\end{center}

\noindent \textbf{Important:} A roll showing a \textbf{10} is never re-rolled, even by Position effects or other mechanics. Critical effects apply only if the overall roll succeeds.

\begin{tcolorbox}[colback=yellow!5!white,colframe=yellow!75!black,title=\textbf{Example: Rolling a 10},breakable]
  Lyra rolls 3 dice: \textbf{10}, 8, 1.  
  She gets: 2 successes (from the 10) + 1 success (from the 8) = 3 successes.  
  The roll also generates 1 SB (from the 1).  
  Because she rolled a 10, the GM grants an extra benefit: ``Your arrow pierces the bandit's shoulder -- he drops his sword.''  
  \textit{(Narrative benefit chosen by GM.)}
\end{tcolorbox}

\section{Boons}
\label{sec:boons}
\index{Boons}

Boons represent momentum and narrative leverage. You may hold up to \textbf{5 Boons}.

\begin{tcolorbox}[colback=yellow!5!white,colframe=yellow!70!black,title=\textit{Lore of the Velvet Touch},breakable]
\emph{``The coin that never spends is the one you don't remember taking.''}  
Boons are the thief's best friend -- a second chance when the lock clicks wrong, a sudden shadow when the guard turns. In Silkstrand, they say a Boon is a debt to luck, and luck always collects.
\end{tcolorbox}

\begin{tcolorbox}[colback=green!5!white,colframe=green!70!black,title=\textbf{Why You Should Spend Boons -- The Math of Fun},breakable]
\index{Boons!spending encouragement}
\textbf{Hoarding Boons for XP conversion leads to less engagement and lower fun scores.}

\medskip
\noindent Every Boon you hold is a dramatic moment waiting to happen:
\begin{itemize}
  \item Re-rolling a critical failure turns despair into triumph.
  \item Improving Position before a risky roll turns a Desperate gamble into a Controlled chance.
  \item Activating an Asset brings your character's network into the story.
\end{itemize}
\textbf{XP from Boons is a trickle (max 2 per session).} The real reward is \textit{agency}. Spend Boons freely -- they are not XP tokens.

\medskip
\noindent \textit{Example:} You have 3 Boons. Instead of converting them later, use 1 to re-roll a missed lockpick, another to improve your Position before a fight, and save the last for an emergency. That is how legends are made.
\end{tcolorbox}

\subsection{Earning Boons}
\label{sec:earning-boons}

\begin{itemize}
  \item \textbf{Partial Success:} Gain 1 Boon (see \ref{sec:outcome-matrix}).
  \item \textbf{Miss:} Gain 2 Boons.
  \item \textbf{Bond-Driven Action:} Once per bond per session, gain 1 Boon when an action meaningfully engages that bond (see Chapter 18 for establishing Bonds).
  \item \textbf{GM Award:} For creativity, sacrifice, or strong narrative play.
\end{itemize}

\noindent \textbf{When not awarded:} Safe rehearsal, null-risk probes, or repeated identical attempts in the same scene do not generate Boons.

\subsection{Spending Boons}
\label{sec:spending-boons}

\begin{itemize}
  \item Re-roll a single die (after seeing the result)
  \item Activate an on-screen Asset (1 Boon)
  \item Improve Position by 1 step before a roll (see \ref{sec:position})
  \item Power certain Rites or abilities (see \ref{ch:magic})
  \item Convert to XP (2 Boons = 1 XP, once per session during downtime, max 2 XP per session)
\end{itemize}

\paragraph{Substituting Fatigue for Boons}
If you have no Boons remaining (or wish to conserve them), you may pay \textbf{1 Fatigue} in place of 1 Boon to activate any effect that normally costs a Boon. This represents pushing through exhaustion, relying on raw effort, or accepting a minor physical toll. The GM may veto this substitution if the fiction makes it implausible (e.g., activating a delicate Asset while exhausted). Using this option does not change the effect's outcome or generate additional Story Beats; it simply changes the resource spent.

\subsection{Carryover Limit}
\label{sec:boon-carryover}
At the end of each scene, reduce held Boons to a maximum of \textbf{2}. Excess are lost -- spend them or lose them.

\section{Position}
\label{sec:position}
\index{Position}

Every action takes place from a \textbf{Position} that reflects the character's advantage.

\begin{tcolorbox}[colback=white!97!gray,colframe=black!80!gray,title=\textbf{Position Effects},breakable]
  \begin{tabularx}{\textwidth}{lXX}
  \toprule
  \textbf{Position} & \textbf{Narrative Frame} & \textbf{Mechanical Effect} \\
  \midrule
  Dominant & You press your advantage & Re-roll one \textit{failure}. \\
  Controlled & Balanced norm & No re-rolls. \\
  Desperate & You act under duress & Re-roll one \textit{success}. \\
  \bottomrule
  \end{tabularx}
  \smallskip
  \noindent You can spend \textbf{1 Boon} to improve Position by one step before rolling.  
  The GM can spend \textbf{2 SB} to worsen your Position by one step (after the roll).
  \end{tcolorbox}

\noindent \textbf{Important:} A roll showing a \textbf{10} is never re-rolled, even by Desperate position.

\begin{tcolorbox}[colback=blue!5!white,colframe=blue!75!black,title=\textit{A Vhasian Duelist's Lesson},breakable]
\emph{``On the Pont-du-Tithe, Position is everything. Stand with the sun at your back -- you are Dominant. Feet on a slippery plank -- Desperate. The bridge does not care for your pride.''}  
Your Position is the story's judgement of your risk. Change the fiction, and Position follows.
\end{tcolorbox}

\paragraph{Boon Interaction with Fatigue:} When Fatigue forces a success re-roll (see \ref{sec:fatigue}), spending 1 Boon negates \textit{one} such re-roll. You may spend Boons after seeing the result of a Fatigue re-roll.

\begin{tcolorbox}[colback=green!5!white,colframe=green!70!black,title=\textbf{Boon Economy at a Glance},breakable]
  \begin{tabular}{ll}
  \toprule
  \textbf{Action} & \textbf{Boon Change} \\
  \midrule
  Partial Success & Gain 1 Boon \\
  Miss & Gain 2 Boons \\
  Bond-Driven Action & Gain 1 Boon (once/bond/session) \\
  GM Award & +1 Boon (rare) \\
  \hline
  Re-roll a die & Spend 1 Boon \\
  Activate Asset & Spend 1 Boon \\
  Improve Position (before roll) & Spend 1 Boon \\
  Convert to XP (downtime) & Spend 2 Boons $\rightarrow$ 1 XP (max 2 XP/session) \\
  \bottomrule
  \end{tabular}
  \medskip
  \textbf{Maximum held:} 5 Boons.  
  At end of scene, reduce to \textbf{2} (excess lost).
\end{tcolorbox}

\begin{effectbox}
    \textbf{Position} (Dominant/Controlled/Desperate) = How risky is this action? -- determines severity of failure and re-rolls.\\
    \textbf{Effect} (Limited/Standard/Great) = How much do you achieve on a success? -- determines scope, scale, or quality.
    
    \vspace{4pt}
    \small
    \textit{Example:} Disarming a trap -- Limited Effect lets one person cross, Standard lets two, Great lets all four.
\end{effectbox}

\section{Story Beats (SB)}
\label{sec:story-beats}
\index{Story Beats}

Story Beats are narrative tools the GM spends to introduce twists. They keep the story alive with complications.

\subsection{SB Spend Examples}
\label{sec:sb-spend}

\begin{itemize}
  \item \textbf{1 SB:} Minor complication -- noise, trace, +1 Supply segment.
  \item \textbf{2 SB:} Moderate setback -- alarm raised, lose cover, lesser foe appears, or advance a faction or front timer by 1 segment.
  \item \textbf{3 SB:} Serious trouble -- reinforcements, key gear breaks.
  \item \textbf{4+ SB:} Major turn -- trap springs, authority arrives, scene shifts.
\end{itemize}

\paragraph{Player Advice:} Don't fear Story Beats -- they are fuel for drama, not punishment. Every time you roll a 1, you give the GM a spark. That spark can light a fire that makes the story burn brighter.

\section{Harm and Fatigue}
\label{sec:harm-fatigue}
\index{Fatigue}\index{Harm}

\subsection{Fatigue Track}
\label{sec:fatigue}
Each character has a Fatigue track equal to \textbf{Body}. Mark Fatigue for exertion, magical strain, or stress.

\begin{itemize}
  \item Each Fatigue step worsens your Position: Dominant $\rightarrow$ Controlled $\rightarrow$ Desperate (see \ref{sec:position}).
  \item If already Desperate, apply --1 die per Fatigue instead.
  \item Penalties from Fatigue and Harm \textbf{stack}. For example, a character with Harm 1 (--1 die) who is also Desperate from Fatigue suffers both penalties.
  \item \textbf{Overflow:} When the track fills, increase Harm by 1 step (see \ref{sec:harm}) and clear all Fatigue.
  \item \textbf{Recovery:} Short rest clears 1--2 Fatigue; full night's rest clears all.
\end{itemize}

\begin{tcolorbox}[colback=brown!5!white,colframe=brown!70!black,title=\textit{Aeler Breath-Counting},breakable]
\emph{``The Deep Watchers say: eight breaths is a corridor, sixteen is a chamber, thirty-two is a fool who forgot to turn back. Fatigue is your breath running out -- mark it, or the mountain will mark you.''}  
Fatigue is not just tiredness; it is the weight of every hard choice you've already made.
\end{tcolorbox}

\subsection{The Roller-Coaster Effect (Taking Harm Clears Fatigue)}
\label{sec:roller-coaster}
\index{Fatigue!roller-coaster effect}
When you take Harm, the shock of the wound temporarily clears your Fatigue. This creates natural dramatic swings -- a character near collapse who takes a wound and suddenly finds a second wind.

\paragraph{Example:} Mateo is at Fatigue 3 (Desperate, --2 dice). A goblin stabs him for Harm 1. Following the procedure in \ref{sec:take-harm-fatigue}, the Harm clears his Fatigue (step 2), then he adds the converted Fatigue (step 3). He ends with Harm 1, Fatigue 0 -- wounded but clear-headed.

\subsection{Harm Levels}
\label{sec:harm}
\begin{center}
\begin{tabular}{ll}
\toprule
\textbf{Harm Level} & \textbf{Effect} \\
\midrule
Harm 1 & --1 die on related actions \\
Harm 2 & --1 die on most actions until treated \\
Harm 3 & Incapacitated or dying \\
\bottomrule
\end{tabular}
\end{center}

\subsection{Taking Harm and Fatigue -- The Procedure (Player Perspective)}
\label{sec:take-harm-fatigue}
\textbf{When you are hit, you report the incoming Harm value and your armor type to the GM. The GM consults the Armor Conversion Table and tells you the result. You then apply the outcome in this order:}

\begin{enumerate}
  \item \textbf{Report to GM:} State the incoming Harm value (e.g., ``I take Harm 2'') and your current armor type (Light/Medium/Heavy).
        \begin{itemize}
          \item \textbf{GM converts using the Armor Conversion Table (see \ref{sec:armor}):}
                \begin{itemize}
                  \item Returns: \textbf{Fatigue gained} ($\geq$ 1 due to minimum rule) and \textbf{Harm remaining} (0 if armor fully converted; >0 if excess Harm bypasses conversion).
                \end{itemize}
          \item \textbf{Example (Player GM Exchange):}
                \begin{itemize}
                  \item \textit{You:} ``I'm hit with Harm 2 while wearing Light Armor.''
                  \item \textit{GM:} ``Consulting table... that's Harm 1 + Fatigue 1.''
                \end{itemize}
        \end{itemize}
  \item \textbf{Apply remaining Harm (if GM says Harm remaining > 0):}
        \begin{itemize}
          \item Clear \textbf{all} existing Fatigue (the shock of injury resets exhaustion).
          \item Increase your Harm level by the Harm remaining value the GM provided.
        \end{itemize}
          \item \textbf{Example:} GM says ``Harm remaining = 1'' $\rightarrow$ Clear Fatigue, Harm +=1.
        \end{itemize}
  \item \textbf{Add the Fatigue} (from GM's conversion result) to your Fatigue track.
        \begin{itemize}
          \item If this fills or exceeds your track (Body value), resolve \textbf{Overflow} (see \ref{sec:fatigue}).
        \end{itemize}
          \item \textbf{Example:} GM says ``Fatigue gained = 2'' $\rightarrow$ Add 2 Fatigue, then check if Fatigue $\geq$ Body.
        \end{itemize}
\end{enumerate}

\begin{tcolorbox}[colback=yellow!5!white,colframe=yellow!75!black,title=\textbf{Player Tip: Trust the GM's Conversion},breakable]
\textbf{Never calculate conversion yourself.} Your job is to:
\begin{itemize}
  \item Accurately report incoming Harm value and armor type.
  \item Trust the GM's returned result (Fatigue gained, Harm remaining).
  \item Mark the resulting Fatigue/Harm changes on your sheet.
  \item Describe the narrative effect (e.g., ``The blow rattles my bones but leaves no wound'' or ``I feel the cut sting as exhaustion vanishes'').
\end{itemize}
\textit{This keeps the pressure engine running: you focus on roleplay, the GM handles mechanics, and the fiction stays vivid.}
\end{tcolorbox}

\paragraph{Armor Conversion Table (Reference for Player-GM Exchange)}
\label{sec:armor}
\begin{center}
\begin{tabular}{l|ccc}
\toprule
\textbf{Armor Type} & \multicolumn{3}{c}{\textbf{Incoming Harm}} \\
\cmidrule(lr){2-4}
& \textbf{Harm 1} & \textbf{Harm 2} & \textbf{Harm 3} \\
\midrule
Light  & Fatigue 1 & Harm 1 + Fatigue 1 & Harm 2 + Fatigue 1 \\
Medium & Fatigue 1 & Fatigue 1 & Harm 2 + Fatigue 1 \\
Heavy  & Fatigue 1 & Fatigue 1 & Fatigue 2 \\
\bottomrule
\end{tabular}
\end{center}
\textit{Minimum 1 Fatigue per hit, even if conversion would give 0.}

\paragraph{Player Examples (Showing GM Interaction)}
\begin{itemize}
  \item \textbf{You report:} ``Harm 2 in Light Armor.''  
        \textit{GM returns:} ``Harm 1 + Fatigue 1.''  
        \textit{You:} Clear existing Fatigue, Harm +=1, then add 1 Fatigue.  
        \textit{Narrative:} ``The blade slides under your shield -- you feel the sting (1 Harm) as exhaustion fades (1 Fatigue added).''
  \item \textbf{You report:} ``Harm 3 in Heavy Armor.''  
        \textit{GM returns:} ``Fatigue 2, Harm remaining 0.''  
        \textit{You:} Add 2 Fatigue (no Harm applied, no Fatigue cleared).  
        \textit{Narrative:} ``Your breastplate absorbs the full impact -- the blow rattles your teeth (2 Fatigue) but draws no blood.''
  \item \textbf{You report:} ``Harm 1 in Heavy Armor.''  
        \textit{GM returns:} ``Fatigue 1, Harm remaining 0.''  
        \textit{You:} Add 1 Fatigue.  
        \textit{Narrative:} ``The glancing blow leaves you winded (1 Fatigue) but unharmed.''
  \item \textbf{Overflow example (Player perspective):}  
        \textit{You (Body 2, current Fatigue 1) report:} ``Harm 3 in Heavy Armor.''  
        \textit{GM returns:} ``Fatigue 2, Harm remaining 0.''  
        \textit{You:} Add 2 Fatigue $\rightarrow$ Fatigue = 3 ($\geq$ Body 2) $\rightarrow$ Overflow: Harm +=1, clear all Fatigue.  
        \textit{Narrative:} ``The crushing blow shatters your guard -- pain flares (Harm 1) as exhaustion vanishes.''
\end{itemize}

\subsection{Recovery}
\label{sec:harm-fatigue-recovery}
\begin{itemize}
  \item \textbf{Fatigue:} Short rest (food/water) removes 2 Fatigue; full night removes all.
  \item \textbf{Harm:} Requires medical care and time -- minor treatment downgrades Harm, proper care removes it.
\end{itemize}

\section{Initiative and Turn Order}
\label{sec:initiative}
\index{Initiative}\index{Turn order}

\paragraph{Default: Narrative Initiative.} Fate's Edge does not use fixed initiative. Turn order follows the fiction:

\begin{itemize}
  \item \textbf{Narrative Fiat:} The GM frames spotlight order based on circumstances.
  \item \textbf{Player Input:} Players may suggest acting when it makes sense.
  \item \textbf{Surprise:} Ambushers act first; targets respond after.
  \item \textbf{Flexibility:} Spotlight may shift mid-scene if fictionally appropriate.
\end{itemize}

\noindent (Optional structured initiative systems -- popcorn, side-based, speed slots -- are available in the GM Guide if your table prefers them.)

\subsection{Turn Economy (Quick Rules)}
\label{sec:turn-economy}
\index{Action economy}\index{Move}\index{Action}

Each character takes \textbf{1 Action and 1 Move} on their turn. Actions and Moves may be taken in any order; repeating the same Action is not allowed unless noted.

\paragraph{Move:}
Traverse up to \textbf{one range band} (Close $\rightarrow$ Near $\rightarrow$ Far, or the reverse).
\begin{itemize}
  \item \textbf{Stand:} Use your Move to stand from prone.
\end{itemize}

\paragraph{Action Options:}

On your action, you may:

\begin{tabular}{|p{2.8cm}|p{6cm}|p{3.5cm}|}
\hline
\textbf{Action} & \textbf{Effect} & \textbf{Notes} \\
\hline
Attack & Make melee or ranged attack against a target. & Ranged in Close is \textbf{Desperate}. \\
\hline
Disengage & Enable safe movement. & Your Move does not provoke opportunity attacks; you cannot Attack this turn. \\
\hline
Shove & Push a target back one range band (Close $\rightarrow$ Near). & Opposed Body + Athletics vs. Body + Athletics. On success, target moves back and suffers --1 Position on their next action. \\
\hline
Grapple & Seize a target to restrain or control. & Opposed Body + Melee vs. Body + Athletics. On success, both are Engaged (cannot move apart). Grapple target suffers --1 die on all actions. Breaking free requires a new opposed roll as an action. \\
\hline
Charge (Full turn) & Move two bands toward an enemy + make one melee attack. & Costs 1 Fatigue, next defense Desperate, requires open ground. \textbf{Leaving Close range during a Charge provokes opportunity attacks as normal.} \\
\hline
Skirmish (Full turn) & Move two bands away from an enemy + make one ranged attack. & Costs 1 Fatigue, next defense Desperate, requires open ground. \textbf{Leaving Close range during a Skirmish provokes opportunity attacks as normal.} \\
\hline
Observe / Change Position & Gain +1 Position once. & Once per turn. \\
\hline
Setup (Teamwork) & Grant an ally +1 Position or step up Effect. & Requires a roll (typically Wits + relevant skill). \\
\hline
Assist (Teamwork) & Spend 1 Boon to give +1 die to an ally's roll. & Only one assistant may use Assist on a given action. Max total assist dice from all sources (including Setup and other abilities) is \textbf{+3}. \\
\hline
Defend / Protect & Adopt a guarding stance. & Your defense rolls improve Position by one step. \\
\hline
Channel / Weave & Runekeeper/ritual preparation. & Interruption may worsen Position. \\
\hline
Cast Rite / Song & Perform a magical action. & May Push (cost Fatigue/Corruption). \\
\hline
Interact & Manipulate an object or environment. & GM sets DV. \\
\hline
Free Items & Short shouts, dropping an item. & Longer actions require proper action types. \\
\hline
\end{tabular}

\paragraph{Sidestep Strike (Closing or Stepping Back)}
\label{sec:sidestep-strike}
\index{Sidestep Strike}\index{Movement during attack}

As part of a melee attack action, you may move \textbf{one range band} (Close $\leftrightarrow$ Near, or Near $\leftrightarrow$ Far) without using your Move action for the turn. If you do so, your Position for that attack worsens by one step (e.g., Controlled becomes Desperate, Dominant becomes Controlled). You cannot also use your regular Move action in the same turn.

\noindent This represents a daring lunge to close distance or a quick step back to create space, sacrificing safety for tactical positioning.

\begin{itemize}
  \item \textbf{Example (Rogue):} Lyra is at Near range with a knife (light weapon). She declares a Sidestep Strike, moving into Close range as she attacks. Her Position drops from Controlled to Desperate, but she now benefits from the knife's Close bonus (+2d). She rolls with the Desperate re-roll (successes are re-rolled once) but gets +2 dice from her weapon.
  \item \textbf{Example (Spear Fighter):} Mateo is at Close range with a spear (medium weapon, +1d at Close, +2d at Near). He steps back to Near as part of his attack, worsening his Position from Controlled to Desperate, but gains +2d instead of +1d.
\end{itemize}

\paragraph{Opportunity Attacks}
\label{sec:opportunity-attacks}
Leaving an enemy's \textbf{Close} range without using \textit{Disengage} provokes one free melee attack (Controlled Position, no action cost). Additionally, the following actions provoke an opportunity attack if performed while \textbf{engaged in Close} with an enemy:
\begin{itemize}
  \item Making a ranged attack (bow, crossbow, thrown weapon).
  \item Casting a spell that does not have the \texttt{[TOUCH]} or \texttt{[WARD]} attribute.
\end{itemize}
A \textbf{Charge} or \textbf{Skirmish} that takes you out of an enemy's Close range provokes an opportunity attack as usual (unless you used Disengage beforehand). Using \textbf{Sidestep Strike} to leave Close range also provokes an opportunity attack unless you also used Disengage.

\paragraph{Position and Dice Modifiers}
\label{sec:position-dice}
\begin{itemize}
  \item \textbf{Position ladder:} Desperate $\rightarrow$ Controlled $\rightarrow$ Dominant.
  \item \textbf{Improving Position:} Each step above Dominant grants \textbf{+1 die} instead.
  \item \textbf{Worsening Position:} Each step below Desperate imposes \textbf{-1 die} instead.
  \item \textbf{Example:} If bonuses would raise Position two steps above Dominant, you gain \textbf{+2 dice}. If penalties would drop you two steps below Desperate, you suffer \textbf{-2 dice}.
\end{itemize}

\paragraph{Caveats and Talents:}
\begin{itemize}
  \item \textbf{Charge} and \textbf{Skirmish} require open ground (no difficult terrain) unless a talent says otherwise (see \ref{ch:talents}).
  \item Talents such as \textit{Effortless Charge} (4 XP) or \textit{Effortless Skirmish} (4 XP) allow ignoring the Fatigue cost once per scene.
  \item The \textit{Relentless} talent (6 XP) allows Charge or Skirmish over difficult terrain (Fatigue cost still applies).
\end{itemize}

\begin{tcolorbox}[colback=gray!5!white,colframe=gray!75!black,title=\textbf{Combat in Brief},breakable]
  \begin{itemize}
    \item \textbf{Turn:} 1 Action + 1 Move.
    \item \textbf{Move:} Close $\leftrightarrow$ Near $\leftrightarrow$ Far.
    \item \textbf{Attack:} Roll Attribute + Skill vs. DV (enemy's defense).  
          Ranged attack while in Close range $\rightarrow$ Desperate.
    \item \textbf{Shove:} Opposed Body + Athletics vs. Body + Athletics. Push target one band.
    \item \textbf{Grapple:} Opposed Body + Melee vs. Body + Athletics. Both Engaged; target suffers --1 die.
    \item \textbf{Sidestep Strike:} Move one band as part of a melee attack; worsen Position by one step.
    \item \textbf{Harm:} When you hit, the target marks Harm (see \ref{sec:take-harm-fatigue}).  
          Enemies use the same Harm levels as players.
    \item \textbf{Fatigue:} Each point worsens your Position (or --1 die if already Desperate).  
          Taking Harm clears Fatigue -- the ``roller-coaster''.
    \item \textbf{Timers:} GMs often use timers for reinforcements, environmental collapse, or enemy morale.
  \end{itemize}
  \textit{Full combat rules are in the(GM Guide).}
\end{tcolorbox}

\section{Weapons \& Armor}
\label{sec:weapons-armor}
\index{Weapons}\index{Armor}\index{Shields}

\begin{tcolorbox}[colback=green!5!white,colframe=green!70!black,title=\textit{Silkstrand Steel},breakable]
\emph{``A Light blade is a courtier's promise -- fast, elegant, easy to hide. A Heavy blade is a condotta captain's ultimatum -- slow, brutal, and final.''}  
Choose your weight as you choose your words. Both can cut.
\end{tcolorbox}

\subsection{Weapons by Weight Class}
\label{sec:weapon-weight}
\begin{itemize}
  \item \textbf{Light (4 XP)} -- fast, concealable.
  \item \textbf{Medium (8 XP)} -- balanced, battlefield standard.
  \item \textbf{Heavy (12 XP)} -- punishing, slow.
\end{itemize}

\subsection{Melee Modifiers}
\label{sec:melee-modifiers}
\begin{center}
\begin{tabular}{llll}
\toprule
\textbf{Weight} & \textbf{Close} & \textbf{Near} & \textbf{Notes} \\
\midrule
Light & +2d & +1d & Quick in tight quarters \\
Medium & +1d & +2d & Set once/scene or --1d first attack \\
Heavy & --1d & +3d & Set once/scene or --2d first attack \\
\bottomrule
\end{tabular}
\end{center}

\subsection{Ranged \& Tempo}
\label{sec:ranged-tempo}
Ranged weapons have a \textbf{Tempo} that determines how they interact with movement, aiming, and reloading.

\paragraph{Tempo Types}
\begin{description}
  \item[Fast] -- You may \textbf{Move and Shoot in the same turn}.  
        \textit{Examples:} throwing knives, light crossbow (if pre-loaded), sling.
  \item[Standard] -- You may either \textbf{Move} or \textbf{Shoot} on your turn, not both.  
        As a special action, you may \textbf{Aim} (gain +1d or +1 Effect on your next shot), but Aim replaces your attack that turn.  
        \textit{Examples:} longbow, heavy crossbow (if pre-loaded), musket.
  \item[Slow] -- The weapon requires a \textbf{full round to load} (or set/brace).  
        \textit{Round 1:} Use your action to load (no attack). You may still move, but moving while loading is at Desperate Position unless you have a talent.  
        \textit{Round 2:} Shoot (and may move afterwards). Then you must reload again for another shot.  
        \textit{Examples:} ballista, hand-cannon, heavy arquebus.
\end{description}

\paragraph{Ranged Modifiers by Weight}
\begin{center}
\begin{tabular}{lllll}
\toprule
\textbf{Weight} & \textbf{Tempo} & \textbf{Close} & \textbf{Near} & \textbf{Far} \\
\midrule
Light (4 XP) & Fast & Controlled & +1d & --- \\
Medium (8 XP) & Standard & Desperate & +2d & +1d \\
Heavy (12 XP) & Slow & Desperate & +1d & +3d \\
\bottomrule
\end{tabular}
\end{center}

\paragraph{Tempo Quick Reference}
\begin{center}
\begin{tabular}{lccc}
\toprule
\textbf{Tempo} & \textbf{Move + Shoot?} & \textbf{Aim Option?} & \textbf{Reload Required} \\
\midrule
Fast      & Yes (same turn) & No      & None (or free) \\
Standard  & No (move or shoot) & Yes (+1d/+1 Effect) & None (Aim costs action) \\
Slow      & No (load one turn, shoot next) & No & Full round \\
\bottomrule
\end{tabular}
\end{center}

\subsection{Weapon Tags (Optional, +4 XP each, max 2)}
\label{sec:weapon-tags}
\index{Weapon tags}
Reach, Close, Accurate, Brutal, Hook, Concealable, Quickdraw, Two-Handed, Off-Hand.

\subsection{Shields (Optional)}
\label{sec:shields}
\begin{center}
\begin{tabular}{llll}
\toprule
\textbf{Shield} & \textbf{XP} & \textbf{Benefit} & \textbf{Tradeoff} \\
\midrule
Buckler & 4 & +1d Defend vs melee or +1 DV once & Off-hand occupied \\
Heater  & 8 & +1d Defend; convert 1 Harm$\rightarrow$1 Fatigue & --1d Ranged \\
Pavise  & 12 & Plant: heavy cover cone & Slow to move; bulky \\
\bottomrule
\end{tabular}
\end{center}

\subsection{Armor}
\label{sec:armor-table}
\begin{center}
\begin{tabular}{llll}
\toprule
\textbf{Armor} & \textbf{XP} & \textbf{Conversion (Harm $\rightarrow$ Fatigue)} & \textbf{Penalty} \\
\midrule
Light  & 4  & 1$\rightarrow$1 (minimum 1 Fatigue per hit) & --- \\
Medium & 8  & 2$\rightarrow$1 (minimum 1 Fatigue per hit) & --1d physical skills \\
Heavy  & 12 & 3$\rightarrow$2 (minimum 1 Fatigue per hit) & --2d physical, no sprint in rough \\
\bottomrule
\end{tabular}
\end{center}
\textit{Notes:} Conversion is applied per Harm instance via GM consultation (see \ref{sec:take-harm-fatigue}). This table shows what the GM will return when you report Harm value and armor type.

\subsection{Condition \& Upkeep}
\label{sec:equipment-condition}
\index{Equipment condition}\index{Neglected}\index{Compromised}
\begin{description}
  \item[Neglected] Weapons: --1d; Armor: conversion worsened by one step (e.g., Light behaves as no armor, Heavy as Medium).
  \item[Compromised] Weapons: --1d first attack each round; Armor: no conversion.
\end{description}
\textit{Fix:} Short rest with tools removes Neglected (no roll required). A scene of work or a smith removes Compromised; a successful Craft roll (DV 3) can reduce the time to a short rest.

\section{Ranged Options}
\label{sec:ranged-options}
\index{Ranged combat}
\begin{itemize}
  \item \textbf{Aim:} +1 die or +1 Effect next shot (lost if you Move or are disrupted).
  \item \textbf{Volley:} Spend extra ammo for +1 die (max +2).
  \item \textbf{Suppress:} Zone under fire; foes there act at --1 die or Limited Effect until they shift cover.
  \item \textbf{Overwatch:} Ready a Controlled shot on a clear trigger crossing your line.
\end{itemize}

\section{Followers and Assets}
\label{sec:followers-assets}
\index{Followers}\index{Assets}

\subsection{Followers}
Followers are named NPCs with a personal relationship to your character. They:
\begin{itemize}
  \item Assist during scenes via Bonds and Assists (see Chapter 18).
  \item Act offscreen to gather information or manage tasks.
  \item Have motives and vulnerabilities -- the GM may place them at risk.
\end{itemize}
Followers are not disposable tools. They can grow, change, or be lost.

\subsubsection{Follower States}
Instead of numeric timers, track a follower's \textbf{Loyalty} and \textbf{Fitness} as one of three states each.

\begin{center}
\small
\begin{tabular}{lll}
\toprule
\textbf{State} & \textbf{Meaning} & \textbf{Mechanical Effect} \\
\midrule
\multicolumn{3}{c}{\textit{Loyalty}} \\
Faithful & Committed, trustworthy & Normal assist, independent actions available \\
Strained & Doubtful, pressured & Assist at --1 die; may refuse risky orders \\
Broken & Betrayed or deserted & Follower leaves or becomes hostile \\
\midrule
\multicolumn{3}{c}{\textit{Fitness}} \\
Ready & Uninjured, rested & Full function \\
Hurt & Wounded, exhausted & Assist at --1 die; no independent actions \\
Down & Incapacitated & Requires rescue or medical care \\
\bottomrule
\end{tabular}
\end{center}

Change states when the fiction demands it (e.g., after a failed mission, a moral compromise, or taking damage). No rolls required.

\subsubsection{High-Capability Followers (Cap 4--5)}
\label{sec:high-cap-followers}
Followers with Capability 4 or 5 are trusted lieutenants, stewards, or commanders. They can be delegated the management of your Assets, freeing you from direct oversight.

\begin{itemize}
  \item A \textbf{Cap 4} follower can be delegated \textbf{up to 2 Assets}.
  \item A \textbf{Cap 5} follower can be delegated \textbf{up to 4 Assets}.
\end{itemize}

When an Asset is delegated to a follower, the follower's own upkeep (paid in XP or downtime action) \textbf{automatically covers the upkeep of those Assets} -- you do not pay separate Asset upkeep. The follower handles the logistics, staff, or political maintenance as part of their role.

If the follower's Loyalty becomes \textbf{Strained} or \textbf{Broken}, all delegated Assets immediately become \textbf{Neglected} as well until the follower is restored. If a follower is \textbf{Compromised}, those assets are lost.

\subsection{Assets}
Assets are durable resources or positions under your control -- safehouses, networks, political offices, relics. They:
\begin{itemize}
  \item Create advantages that persist beyond a single scene.
  \item Resolve problems that dice alone cannot.
  \item Cost Boons to activate (the GM determines the impact).
\end{itemize}
Assets can be threatened, degraded, or lost. Overuse attracts attention.

\subsubsection{Asset Status}
Assets have a single unified \textbf{Status} that describes their overall health and usability.

\begin{center}
\small
\begin{tabular}{lll}
\toprule
\textbf{Status} & \textbf{Meaning} & \textbf{Effect} \\
\midrule
Flourishing & Well-maintained, respected & +1 die on rolls involving the asset, +1 Yield \\
Stable & Normal operation & No modifier \\
Strained & Problems accumulating & --1 die on asset-related rolls, no Yield \\
Collapsed & Non-functional & Requires quest or major downtime to restore \\
\bottomrule
\end{tabular}
\end{center}

The Status worsens when the asset is neglected, attacked, or suffers a public failure. It improves through dedicated downtime actions or successful missions that benefit the asset.

\subsection{Followers vs. Assets}
\begin{center}
\begin{tabular}{p{3cm}p{5cm}p{5cm}}
\toprule
 & \textbf{Followers} & \textbf{Assets} \\
\midrule
Nature & People & Resources/Structures \\
Primary Role & Personal support & Strategic leverage \\
Activation & Assists (spend Boon) or narrative interaction & Boon expenditure \\
Risk & Emotional, narrative & Structural, political \\
Growth & Can become PCs & Can expand or collapse \\
\bottomrule
\end{tabular}
\end{center}

\paragraph{Design Principle:} Followers create \textit{depth}. Assets create \textit{reach}. High-tier play depends on managing both.

\section{Game Structure}
\label{sec:game-structure}
\index{Game structure}\index{Time scales}

\subsection{Time Scales}
\begin{description}
  \item[Moment] A heartbeat, a single action.
  \item[Some Time] A few minutes, a short activity.
  \item[Significant Time] Hours, extended effort.
  \item[Days] Large-scale endeavors.
\end{description}

\subsection{Game Units}
\begin{description}
  \item[Scene] Basic narrative unit, covers a specific conflict.
  \item[Session] One game session (3--6 hours).
  \item[Campaign] Entire story arc.
\end{description}

\section{Action Resolution Steps}
\label{sec:action-resolution}
\index{Action resolution}

\begin{enumerate}
  \item Describe your intent and method.
  \item Build dice pool: Attribute + Skill (+ gear, assists).
  \item Roll d10s; count Successes (6+) and Story Beats (1s).
  \item Compare Successes to DV.
  \item Apply outcome from the Outcome Matrix (see \ref{sec:outcome-matrix}).
  \item GM spends SB if applicable (see \ref{sec:story-beats}).
  \item Earn Boons for Partial or Miss (see \ref{sec:boons}).
\end{enumerate}

\begin{tcolorbox}[colback=blue!5!white,colframe=blue!75!black,title=\textbf{Quick Reference},fonttitle=\bfseries,breakable]
\textbf{Dice Pool:} Attribute + Skill d10s \\
\textbf{Success:} Die $\geq$ 6 \\
\textbf{Story Beat (SB):} Die = 1 gives SB to GM \\
\textbf{DV:} 2 (easy) to 5+ (extreme) -- default 3 \\
\textbf{Harm:} 3-level system; see \ref{sec:take-harm-fatigue} for procedure \\
\textbf{Boons:} 2 on miss, 1 on partial \\
\textbf{Armor Conversion:} Report Harm value + armor type to GM; receive Fatigue gained + Harm remaining \\
\textbf{Position:} Dominant (re-roll failure), Controlled (normal), Desperate (re-roll success) \\
\textbf{Fatigue:} Worsens Position (or --1 die if Desperate). Full track $\rightarrow$ Harm+1, clear Fatigue. \\
\textbf{Harm:} 1 (--1 die), 2 (--2 dice), 3 (incapacitated). Taking Harm clears Fatigue first. \\
\end{tcolorbox}

\section{Narrative Suggestions}
\label{sec:narrative-suggestions}
\index{Narrative play}\index{Flavor}

\begin{itemize}
  \item \textbf{Collaborative Scene Framing:} Players may suggest scene elements (weather, NPC reactions) with GM approval.
  \item \textbf{Intent-Driven Resolution:} If success is reasonably assured, the GM may ask players to describe \textit{how} rather than roll.
  \item \textbf{Flashback Declarations:} Spend 1 Boon to declare a flashback establishing past preparation, and describe the scene.
  \item \textbf{Descriptive Assistance:} Players can assist by providing vivid descriptions, granting a +1 die to the primary actor.
  \item \textbf{Proactive Storytelling:} Players may suggest minor favorable details (a helpful NPC, a useful item on hand) subject to GM approval.
\end{itemize}

\begin{tcolorbox}[colback=black!5!white,colframe=black!60!white,title=\textbf{Core Mechanics Quick Reference},breakable]
  \begin{center}
  \small
  \begin{tabular}{ll}
  \toprule
  \textbf{What} & \textbf{How} \\
  \midrule
  Dice Pool & Attribute + Skill d10s \\
  Success & Die $\geq$ 6 \\
  Story Beat (SB) & Die = 1 \\
  Difficulty (DV) & Default 3 (2=easy, 4=hard, 5+=extreme) \\
  Outcomes & Clean Success (S$\geq$DV, SB=0) \\ & Success with SB (S$\geq$DV, SB>0) \\ & Partial (0<S<DV $\rightarrow$ gain 1 Boon) \\ & Miss (S=0 $\rightarrow$ gain 2 Boons) \\
  Boons & Max 5; trim to 2 at scene end. Spend: re-roll, improve Position, activate Asset, convert 2$\rightarrow$1 XP. \\
  Position & Dominant (re-roll failure), Controlled (normal), Desperate (re-roll success) \\
  Fatigue & Worsens Position (or --1 die if Desperate). Full track $\rightarrow$ Harm+1, clear Fatigue. \\
  Harm & 1 (--1 die), 2 (--2 dice), 3 (incapacitated). Taking Harm clears Fatigue first. \\
  Armor Conversion & Report Harm value + armor type to GM; apply returned Fatigue gained/Harm remaining \\
  Sidestep Strike & Move one band as part of melee attack; worsen Position by one step. \\
  Shove & Opposed Body + Athletics. Push target one band. \\
  Grapple & Opposed Body + Melee. Both Engaged; target suffers --1 die. \\
  \bottomrule
  \end{tabular}
  \end{center}
\end{tcolorbox}
